\section{结论}
\label{sec:conclusion}
%本文引入了用于验证离散时间动力系统 LTL 性质的转移证书。与状态三元组方法~\cite{wongpiromsarn2015automata} 相比，转移证书降低了保守性，并且其搜索空间不大于 $\mathcal{X} \times Q \times Q$，显著小于闭包证书~\cite{murali2024closure}。我们给出了合成方法，并在多个案例中展示了相对于闭包证书的加速。未来工作将重点研究其在控制器合成中的应用，并将这些概念适配到连续时间系统的验证与合成中。

本文提出了转移证书，作为一种用于验证具有连续状态空间的离散时间动力系统线性时序逻辑（LTL）性质的高效框架。我们的方法利用转移安全性证书和转移持久性证书，相比状态三元组方法~\cite{wongpiromsarn2015automata} 提供了更大的灵活性，并相对于闭包证书~\cite{murali2024closure} 减少了搜索空间，从而在 LTL 验证中实现显著加速。实验结果表明，在因接受状态可达而导致状态三元组方法结构上不适用的基准上，我们的转移持久性证书能够成功验证该性质；而在状态三元组方法成功的基准上，我们的方法在提供额外灵活性的同时实现了可比较的性能。我们开发了基于多项式模板和神经网络模板的自动化合成方法，并通过 Kuramoto 振子和双房间温度模型等案例研究展示了其实用性。尽管如此，我们的方法仍存在若干不足，这些不足也启发了未来工作。第一，对于具有高度非线性动力学的高维系统，合成过程，尤其是神经网络模板的合成，可能由于 SMT 求解瓶颈而在计算上较为昂贵。第二，从 LTL 公式构造 NBA 在最坏情况下相对于公式大小具有指数复杂度；然而，Spot~\cite{spot} 和 Owl~\cite{owl} 等现代转换器对于大多数实际规约都能产生紧凑自动机，因此这在实践中很少成为瓶颈。此外，当前工作关注离散时间系统，而连续时间系统仍是一个开放挑战。
